\chapter{Introduction}\label{intro}
\section{Dissertation Outline}

In the digital era, recommender systems (RecSys) play a crucial role in helping users discover relevant items across online platforms. By modeling historical user–item interactions, RecSys learns user preferences and provides personalized recommendations. To effectively capture collaborative signals, graph learning has become a dominant paradigm in RecSys, where users and items are organized into a bipartite structure for information propagation~\cite{glrssurvey21}. However, user–item interactions are inherently sparse due to limited user exploration and the extremely large candidate item space. Under such sparse settings, relying solely on pairwise interaction modeling remains insufficient for robust personalization. To address this limitation, this thesis explores the integration of complementary information sources into recommendation.
Inspired by the idea of leveraging multiple forms of knowledge simultaneously, we investigate two primary directions, as well as their integration, to enhance personalized recommendation.

First, we introduce additional relations and attributes into graph learning. Beyond conventional bipartite graphs that model only pairwise user–item interactions, we study hypergraph learning to capture higher-order relations and complex dependencies among multiple entities. Hypergraphs connect multiple nodes within a single hyperedge, allowing the incorporation of side information and structured relations beyond direct interactions. By modeling richer structural patterns, we enhance representation learning and alleviate sparsity issues in RecSys.

Second, we incorporate prior knowledge from LLMs into recommendation. Traditional embedding-based graph models rely on fixed-length vector representations, which may fail to preserve the full semantic richness of textual descriptions and world knowledge. Recent advances in LLMs demonstrate strong capabilities in understanding semantic relations and functional roles of items derived from large-scale text corpora. In this thesis, we leverage LLMs to inject semantic priors and reasoning ability into RecSys, while retaining graph models as efficient tools for interaction modeling. Graph learning captures collaborative and structural patterns, whereas LLM reasoning contributes high-level semantic understanding, forming a complementary framework for personalization.

This thesis explores the challenges and opportunities of integrating graph learning and LLM reasoning to enhance RecSys. Compared with direct user–item interactions that provide explicit preference signals, auxiliary relations, structured side information, and LLM-based prior knowledge can be heterogeneous, multi-sourced, and partially misaligned with recommendation objectives. By addressing these challenges properly, we demonstrate significant improvements in recommendation performance. The detailed introduction for the work presented is as follows.

\section{Hypergraph Learning for Relation Modeling in Recommendation}
(Part of this section was previously published in~\cite{uprth}.)

Although graph-based recommender systems have shown strong performance, they mainly model pairwise user–item interactions, which limits their ability to capture higher-order relations and heterogeneous signals. In real-world recommendation scenarios, rich relational structures naturally exist, including multi-entity dependencies and auxiliary supervision signals. Hypergraph learning provides a principled framework to model such complex relations, as a hyperedge can connect multiple nodes simultaneously.

In this work, we propose a multitask framework named Unified Pretraining for Recommendation via Task Hypergraphs (UPRTH). To provide a unified learning paradigm for diverse relational signals, we generalize different pretext tasks as hyperedge prediction problems within a task hypergraph. In this way, heterogeneous relations and auxiliary attributes can be integrated into recommendation through a shared hypergraph structure. Furthermore, a transitional attention layer is designed to adaptively learn the relevance between each auxiliary task and the main recommendation objective. Experimental results on three benchmark datasets verify the effectiveness of hypergraph learning for modeling diverse relations in recommendation.

\section{Hypergraph Pretraining Guided by Text Prompt with PLM}
(Part of this section was previously published in~\cite{ihp}.)

Building upon the flexibility of hypergraph construction introduced in UPRTH, we extend this paradigm to cross-domain recommendation by incorporating semantic guidance into hypergraph pretraining. In cross-domain settings where domains share the same users but completely different items, discrepancy in data distribution significantly hinder effective knowledge transfer via structural signals alone. To address this limitation, we introduce plain-text prompts that explicitly describe task semantics and domain connections, and encode them using a pretrained LLM to guide hypergraph pretraining. To be concrete, we propose Instruction-based Hypergraph Pretraining (IHP), a framework that injects LLM-encoded textual instructions into hypergraph learning to align source-domain knowledge with target-domain recommendation objectives. A prompt hypergraph convolution layer is designed to inject the encoded prompt representations into hyperedge-level information propagation, enabling instruction-aware high-order relation modeling. By integrating semantic guidance from text prompts with structural hypergraph learning, IHP enhances cross-domain knowledge transfer and improves recommendation performance across multiple datasets.

\section{Training Large Recommendation Models via Graph-LLM Token Alignment}
(Part of this section was previously published in~\cite{glta}.)

While IHP introduces plain-text prompts to provide task-level semantic guidance for hypergraph learning, the amount of injected knowledge remains limited to a small number of instructions. To further integrate richer prior knowledge from LLMs, we explore a deeper alignment between graph-based recommendation models and pretrained language representations. Traditional collaborative filtering models primarily rely on user–item interactions and lack direct access to the extensive semantic knowledge encoded in LLMs. At the same time, directly applying LLMs for recommendation often suffers from ambiguous item generation and poor scalability.
To address these challenges, we propose training large recommendation models via Graph-LLM Token Alignment (GLTA). Instead of generating items through free-form text, GLTA aligns user and item representations from the interaction graph with pretrained LLM token embeddings, allowing the graph model to inherit semantic knowledge from frozen LLMs in a structured manner. Furthermore, we introduce Graph-Language Logits Matching (GLLM) to optimize token alignment for end-to-end item prediction, eliminating ambiguity in text generation while preserving scalability. Experimental results on benchmark datasets demonstrate that token-level alignment effectively enhances recommendation performance.

\section{LLM-Inferred Substitute and Complement Signals for Recommendation}
As the size of LLMs continues to grow, integrating them directly into graph model training becomes increasingly computationally expensive. To achieve a more lightweight yet effective integration of LLM knowledge and generative ability into recommendation, we instead explore leveraging the world knowledge encoded in LLMs to evaluate functional relations between items and enhance recommendation through refined relational signals.

Substitution and complementarity are essential item–item relations for understanding user intent in recommendation systems. However, reliable annotations of such relations are rarely available in public datasets, and behavior-based signals such as co-view and co-purchase often conflate true semantic relations with incidental co-occurrences. To address this challenge, we propose Enhancing Recommendation with LLM-Inferred Substitute and Complement Signals (ERLSC), a framework that leverages LLMs to assess functional relations within item groups and guide hyperedge weighting. We adopt a hypergraph formulation to represent item groups induced by co-occurrence signals, and use LLMs to evaluate pairwise substitute and complementary relations within each group. Based on these semantic assessments, we design an LLM-Guided Hypergraph Convolution layer that initializes and adaptively updates hyperedge weights, enabling the model to emphasize semantically meaningful item groups while suppressing noisy co-occurrences. Extensive experiments on three public datasets demonstrate that ERLSC consistently outperforms strong baselines, and ablation studies verify the effectiveness of each component.

\section{Dynamic Recommendation with Item-Item relations via LLM-based Agents}
(Part of this section was previously published in~\cite{agentdr}.)

While ERLSC leverages LLMs to refine structural signals through LLM-guided hyperedge weighting within a unified graph model, it still follows a single-model learning paradigm. In this setting, the final recommendation performance ultimately relies on a single trained model, even if enhanced by LLM-derived relational signals. Motivated by the strong reasoning capabilities of modern LLMs, we instead explore an LLM-based agent framework that can coordinate multiple recommendation tools and perform explicit relational reasoning at inference time, without being restricted to the output of a single recommendation model.

Recent agent-based recommendation methods use memory and prompting to mimic user behavior, but they often hallucinate non-existent items and struggle with full-catalog ranking. To address these challenges, we propose AgentDR, an LLM-agent framework that combines commonsense reasoning with scalable recommendation tools. AgentDR treats the LLM as a high-level decision maker: full-catalog ranking is delegated to traditional models for efficiency and ID grounding, while the LLM integrates multiple recommendation outputs based on personalized tool suitability. Building on the effectiveness of leveraging LLM world knowledge to evaluate item–item relations in ERLSC, we further enable the agent to explicitly reason over substitute and complementary relations inferred from the user’s historical interactions, guiding the selection and refinement of candidate items. This modular design preserves scalability and stability while adding relational reasoning at the decision level. Experiments on three public grocery datasets show that AgentDR substantially improves full-ranking performance, achieving on average a twofold gain over its underlying tools, and we further introduce an LLM-based evaluation metric that jointly measures semantic alignment and ranking correctness.